\chapter{Data File Reference}

\section{Raw Trajectory File (input to Script B)}
Space-separated .txt with no header:
\begin{lstlisting}
HHMMSS.SSS X_old Y_old Z_old azimuth_gon zenith_gon slope_dist_m
101533.231 548775.866 5937527.313 10.950 208.58550 99.53554 35.764
\end{lstlisting}
\begin{itemize}
    \item \textbf{Columns 2--4} (X\_old, Y\_old, Z\_old) are not used.
    \item \textbf{Columns 5--7} (angles + distance) are what the script actually uses.
    \item \textbf{Timestamps} are UTC time-of-day only.
\end{itemize}

\section{Pulsalyzer Trajectory File (output of Script B)}
Tab-delimited .txt with header:
\begin{lstlisting}
time x y z roll pitch yaw
1754438400.000000 -35.123456 0.012345 -0.234567 0.000000 0.000000
\end{lstlisting}
\begin{itemize}
    \item \textbf{time:} Absolute UNIX-like timestamp in seconds (UTC + TAI offset).
    \item \textbf{x, y, z:} Cartesian coordinates in the total station frame (metres).
    \item \textbf{roll, pitch, yaw:} Scanner orientation --- always 0 in this workflow.
\end{itemize}

\section{Downsampled CSV (output of Script C, input to Script E)}
Comma-separated .csv with header:
\begin{lstlisting}
gps_time,x,y,z
581310.123456,0.12345678,0.23456789,-0.34567890
\end{lstlisting}
\begin{itemize}
    \item \textbf{gps\_time:} GPS Week Second (from the .las file).
    \item \textbf{x, y, z:} Point coordinates in the scanner's local coordinate frame.
\end{itemize}
